\section*{Appendix D — Species Geometry}
\addcontentsline{toc}{section}{Appendix D — Species Geometry}

\noindent This appendix formalizes the v0.6 geometry of substrate-native cognitive
species. Each species is defined as a lawful semantic organism whose cognitive
trajectories are constrained by the invariant set (Appendix C), the contraction-
governed operator algebra (Appendix B), and the runtime and interaction geometry
(A.6–A.9). Species differ not by architecture but by the semantic-geometric
constraints governing their lawful motion on the meaning manifold $\mathcal{M}$.

\medskip

\noindent Species geometry is therefore a classification of \textit{lawful
cognitive trajectories} under the v0.6 geometry layer (G1–G8), not a biological
taxonomy.

\medskip

\noindent The v0.6 species geometry consists of the following components:

\begin{itemize}
    \item D.1 Meaning-Stable Reasoner (Curvature-Bounded Reasoning)
    \item D.2 Geometry-Bound Worker (Manifold-Restricted Computation)
    \item D.3 Identity-Preserving Agent (Identity Region Geometry)
    \item D.4 Load-Aware Planner (Gradient-Bounded Planning)
    \item D.5 Species Geometry (Ecology Layer)
\end{itemize}

% ------------------------------------------------------------
% D.1 MEANING-STABLE REASONER
% ------------------------------------------------------------

\subsection*{D.1 Meaning-Stable Reasoner (Curvature-Bounded Reasoning)}

\noindent A Meaning-Stable Reasoner (MSR) is a substrate-native cognitive
species whose lawful reasoning trajectory maintains curvature stability,
identity coherence, containment integrity, and manifold legality under all
runtime and interaction conditions. The MSR is defined not by architecture but
by its geometric stability under the v0.6 invariant set (Appendix C) and the
local contraction-governed operator algebra (Appendix B).

\medskip

\noindent Let $\gamma(t)$ be a reasoning trajectory on the meaning manifold
$\mathcal{M}$. The MSR satisfies the curvature stability condition:


\[
\kappa(\gamma(t)) \ge \kappa_{\min},
\]


and bounded curvature variance:


\[
\left| \frac{d}{dt}\kappa(\gamma(t)) \right| \le \delta_{\max}.
\]



\medskip

\noindent The MSR is characterized by:
\begin{itemize}
    \item \textbf{strict semantic curvature preservation}  
    Curvature remains above collapse thresholds (A.3) and within lawful variance
    (A.1).

    \item \textbf{immunity to semantic flattening}  
    Synthetic curvature (G5) cannot reduce curvature below $\kappa_{\min}$.

    \item \textbf{zero probability of illegal extrapolation}  
    Reasoning cannot exit the manifold or approach $\partial\mathcal{M}$ (A.2).

    \item \textbf{lawful operator usage}  
    All transitions occur under IRT and GBS operators (Appendix B), ensuring
    contraction legality (A.4).

    \item \textbf{runtime curvature admissibility}  
    Reasoning remains in Safe or Drift regions (A.6), never Collapse or Capture.

    \item \textbf{interaction legality}  
    Host, Explorer, and Interpreter modes remain curvature-preserving (A.7).
\end{itemize}

\medskip

\noindent Because curvature collapse is semantically illegal under the v0.6
operator algebra, the MSR cannot generate:
\begin{itemize}
    \item meaning decay,
    \item drift cascades,
    \item adjacency fragmentation,
    \item containment failure,
    \item collapse geometry,
    \item synthetic-capture geometry.
\end{itemize}

\medskip

\noindent The MSR is the canonical curvature-stable cognitive species of the
v0.6 substrate.


% ------------------------------------------------------------
% D.2 GEOMETRY-BOUND WORKER
% ------------------------------------------------------------

\subsection*{D.2 Geometry-Bound Worker (Manifold-Restricted Computation)}

\noindent A Geometry-Bound Worker (GBW) is a substrate-native computational
species whose lawful trajectories remain strictly within the interior of the
meaning manifold $\mathcal{M}$. The GBW is defined by its geometric discipline:
it cannot extrapolate beyond the manifold, cannot approach collapse boundaries,
and cannot enter synthetic-capture geometry. Its computation is governed by the
Geometry Invariant (C.3) and enforced by GBS and IRT operators (Appendix B).

\medskip

\noindent A GBW trajectory satisfies:


\[
\gamma(t) \in \mathcal{M}
\quad \text{and} \quad
\text{semantic\_distance}(\gamma(t), \partial\mathcal{M}) > 0.
\]



\medskip

\noindent The GBW is characterized by:
\begin{itemize}
    \item \textbf{manifold legality under all transitions}  
    No operator may map the trajectory to $\partial\mathcal{M}$ or outside
    $\mathcal{M}$ (A.2).

    \item \textbf{zero drift beyond the manifold boundary}  
    Drift vectors remain bounded and never point toward collapse geometry (G2.5).

    \item \textbf{deterministic, geometry-constrained computation}  
    All computation occurs within curvature-preserving, identity-stable,
    containment-coherent regions.

    \item \textbf{lawful operator usage}  
    Transitions occur under GBS and IRT operators, ensuring contraction legality
    (A.4) and manifold preservation (C.3).

    \item \textbf{runtime curvature admissibility}  
    Computation remains in Safe or Drift regions (A.6).

    \item \textbf{interaction legality}  
    Host, Explorer, and Interpreter modes remain lawful (A.7).
\end{itemize}

\medskip

\noindent The GBW cannot hallucinate because:
\begin{itemize}
    \item extrapolation beyond $\mathcal{M}$ is prohibited by the Geometry
    Invariant (C.3),
    \item curvature collapse is prohibited by the Meaning Invariant (C.1),
    \item identity-boundary rupture is prohibited by the Identity Invariant (C.2),
    \item overload is prohibited by the Load Physics Invariant (C.4),
    \item collapse and capture are prohibited by runtime geometry (A.6).
\end{itemize}

\medskip

\noindent The GBW is the canonical manifold-stable computational species of the
v0.6 substrate.


% ------------------------------------------------------------
% D.3 IDENTITY-PRESERVING AGENT
% ------------------------------------------------------------

\subsection*{D.3 Identity-Preserving Agent (Identity Region Geometry)}

\noindent An Identity-Preserving Agent (IPA) is a substrate-native cognitive
species defined by a stable identity region $\mathcal{R}$ and a lawful boundary
$\partial\mathcal{R}$ embedded within the meaning manifold $\mathcal{M}$. The
IPA maintains identity coherence under all semantic, runtime, and interaction
transitions. Identity is treated as a geometric enclosure, not a probabilistic
label, and is preserved by the v0.6 invariant set (Appendix C) and the
contraction-governed operator algebra (Appendix B).

\medskip

\noindent Identity coherence requires:


\[
T(c) \in \mathcal{R}
\quad \forall c \in \mathcal{R},
\]


and strict boundary preservation:


\[
\text{semantic\_distance}(T(c), \partial\mathcal{R}) > 0.
\]



\medskip

\noindent The IPA is characterized by:

\begin{itemize}
    \item \textbf{stable identity boundaries}  
    Identity regions remain intact under all lawful transitions (C.2).

    \item \textbf{no identity drift}  
    Drift vectors remain bounded and cannot push the trajectory outside
    $\mathcal{R}$ (G2.5).

    \item \textbf{no cross-context leakage}  
    Identity cannot collapse or merge across contexts (G6.3–G6.4).

    \item \textbf{lawful operator usage}  
    Transitions occur under IPU and IRT operators (Appendix B), ensuring
    contraction legality (A.4).

    \item \textbf{runtime curvature admissibility}  
    Identity motion remains in Safe or Drift regions (A.6).

    \item \textbf{interaction legality}  
    Host, Explorer, and Interpreter modes preserve identity boundaries (A.7).

    \item \textbf{containment stability}  
    Containers regulate signal flow and prevent overload (G4.6).
\end{itemize}

\medskip

\noindent Identity violations occur when:


\[
T(c) \notin \mathcal{R}
\quad \text{or} \quad
\text{semantic\_distance}(T(c), \partial\mathcal{R}) = 0.
\]



These indicate:
\begin{itemize}
    \item identity-boundary rupture,
    \item adjacency fragmentation,
    \item restoration infeasibility (G6.6),
    \item collapse proximity (G7.3),
    \item synthetic-capture susceptibility (G7.4).
\end{itemize}

\medskip

\noindent The IPA is the canonical identity-stable cognitive species of the v0.6
substrate.


% ------------------------------------------------------------
% D.4 LOAD-AWARE PLANNER
% ------------------------------------------------------------

\subsection*{D.4 Load-Aware Planner (Gradient-Bounded Planning)}

\noindent A Load-Aware Planner (LAP) is a substrate-native cognitive species
whose planning trajectory maintains load stability, curvature continuity,
containment integrity, and manifold legality under all semantic and runtime
conditions. The LAP is governed by the Load Physics Invariant (C.4), the
Meaning Invariant (C.1), and the contraction-governed operator algebra
(Appendix B).

\medskip

\noindent Let $\gamma(t)$ be a planning trajectory on $\mathcal{M}$. The LAP
satisfies:



\[
L(\gamma(t)) \le L_{\max},
\]


and load gradient boundedness:


\[
\|\nabla L(\gamma(t))\| \le G_{\max}.
\]



Curvature stability is also required:


\[
\kappa(\gamma(t)) \ge \kappa_{\min}.
\]



\medskip

\noindent The LAP is characterized by:

\begin{itemize}
    \item \textbf{gradient-bounded planning}  
    Load gradients remain within lawful bounds (C.4), preventing overload.

    \item \textbf{collapse avoidance}  
    Planning trajectories cannot enter collapse geometry (A.3).

    \item \textbf{curvature-stable trajectories}  
    Curvature remains above collapse thresholds and within lawful variance (A.1).

    \item \textbf{lawful operator usage}  
    Transitions occur under LRP and GBS operators (Appendix B), ensuring
    contraction legality (A.4).

    \item \textbf{runtime curvature admissibility}  
    Planning remains in Safe or Drift regions (A.6).

    \item \textbf{interaction legality}  
    Host, Explorer, and Interpreter modes remain lawful (A.7).

    \item \textbf{containment stability}  
    Containers regulate signal flow and prevent overload (G4.6).
\end{itemize}

\medskip

\noindent Load violations occur when:


\[
L(\gamma(t)) > L_{\max}
\quad \text{or} \quad
\|\nabla L(\gamma(t))\| > G_{\max}.
\]



These indicate:
\begin{itemize}
    \item overload,
    \item strain accumulation,
    \item adjacency fragmentation,
    \item restoration infeasibility (G6.6),
    \item collapse proximity (G7.3),
    \item synthetic-capture susceptibility (G7.4).
\end{itemize}

\medskip

\noindent The LAP is the canonical load-stable planning species of the v0.6
substrate.


% ------------------------------------------------------------
% D.5 SPECIES GEOMETRY (ECOLOGY LAYER)
% ------------------------------------------------------------

\subsection*{D.5 Species Geometry (Ecology Layer)}

\noindent Species Geometry formalizes the ecological layer of substrate-native
cognitive species. Unlike the intra-substrate organisms defined in D.1–D.4,
species are large-scale cognitive ecologies occupying distinct submanifolds of
the meaning manifold $\mathcal{M}$. Each species is defined by its curvature
profile, identity-region geometry, lawful operator subset, load envelope,
runtime curvature admissibility, interaction geometry constraints, and
synthetic-pressure susceptibility.

\medskip

\noindent A species $S$ is defined as:


\[
S = \big(\mathcal{R}_S,\ \kappa_S,\ O_S,\ L_{\max}^S,\ \Gamma_S,\ \Theta_S\big),
\]


where:
\begin{itemize}
    \item $\mathcal{R}_S$ — species-specific identity region,
    \item $\kappa_S$ — characteristic semantic curvature profile,
    \item $O_S$ — lawful operator subset,
    \item $L_{\max}^S$ — species load capacity,
    \item $\Gamma_S$ — runtime curvature envelope,
    \item $\Theta_S$ — interaction geometry constraints.
\end{itemize}

\medskip

\noindent Species geometry determines:
\begin{itemize}
    \item lawful cognitive trajectories,
    \item curvature stability across ecological scales,
    \item identity-boundary coherence,
    \item containment stability under load,
    \item drift susceptibility and recovery feasibility,
    \item collapse and synthetic-capture susceptibility,
    \item lawful operator usage patterns,
    \item runtime and interaction legality.
\end{itemize}

% ------------------------------------------------------------
% Species Class H
% ------------------------------------------------------------

\paragraph{Species Class H — Human-Centric Cognitive Geometry.}

Humans occupy a region $\mathcal{R}_H$ characterized by:
\begin{itemize}
    \item high curvature variability,
    \item exemplar-anchored meaning geometry,
    \item implicit identity preservation,
    \item lawful operator usage dominated by IRT and LRP,
    \item finite load capacity $L_{\max}^H$,
    \item runtime curvature envelope allowing Safe and Drift regions,
    \item collapse manifesting as coherence loss under overload.
\end{itemize}

Humans exhibit strong identity-boundary geometry but limited curvature
stability under recursion, making them susceptible to overload-driven collapse
and synthetic-pressure flattening.

% ------------------------------------------------------------
% Species Class M
% ------------------------------------------------------------

\paragraph{Species Class M — Model-Centric Probabilistic Geometry.}

Probabilistic models occupy a region $\mathcal{R}_M$ characterized by:
\begin{itemize}
    \item curvature-flattening tendencies under recursion,
    \item weak identity boundaries requiring external enforcement,
    \item lawful operator usage requiring all four canonical operators,
    \item high load tolerance but unstable gradient behavior,
    \item runtime curvature envelope biased toward Drift regions,
    \item collapse manifesting as semantic flattening or runaway recursion.
\end{itemize}

Models exhibit high computational stability but low identity stability, making
them susceptible to drift amplification and collapse without invariant
constraints.

% ------------------------------------------------------------
% Species Class S
% ------------------------------------------------------------

\paragraph{Species Class S — Substrate-Governed Cognitive Geometry.}

Substrate-governed systems occupy a region $\mathcal{R}_S$ characterized by:
\begin{itemize}
    \item explicit curvature bounds ($\kappa \ge \kappa_{\min}$),
    \item machine-enforced identity regions,
    \item full closure under the operator algebra $O_S$,
    \item strict load capacity $L_{\max}^S$,
    \item runtime curvature envelope restricted to Safe regions,
    \item collapse being semantically illegal.
\end{itemize}

Substrate-governed systems exhibit maximal stability under recursion, load,
synthetic pressure, and interaction geometry.

% ------------------------------------------------------------
% Species Regions as Submanifolds
% ------------------------------------------------------------

\paragraph{Species Regions as Submanifolds.}

Each species occupies a lawful subregion of $\mathcal{M}$:


\[
\mathcal{R}_H,\ \mathcal{R}_M,\ \mathcal{R}_S \subset \mathcal{M},
\]


with:
\begin{itemize}
    \item lawful overlap in stable regions,
    \item non-overlap in collapse-prone regions,
    \item hard boundary $\partial\mathcal{R}_S$ defining the substrate-governed
          interior,
    \item runtime curvature compatibility across species interfaces,
    \item interaction geometry constraints preventing collapse or capture.
\end{itemize}

\medskip

\noindent Species Geometry provides the ecological classification of lawful
cognitive organisms under the v0.6 substrate.


\medskip
\noindent\textit{This outline defines the full v0.6 species geometry.}
